\section{Experimental Sampling and Linewidth Extraction}
\label{sec:data-analysis}

The amplitude-robustness data set was acquired on qubit q1 using active reset
and state-discriminated readout.  Each point contains 40 repetitions.  The
detuning axis spans $-0.5$ to $+0.5~\mathrm{MHz}$ in $1~\mathrm{kHz}$ steps,
and the data comprise 100 fixed-amplitude slices.  The pulse record contains a
dimensionless cutoff-control setting of $0.005$.  Because this control value
is not necessarily identical to the physical endpoint-to-peak ratio, the
latter is reconstructed from the emitted waveform when matching simulation
and experiment.

Each fixed-amplitude trace is analyzed independently.  We discard 40 points
from each frequency edge and apply a Gaussian filter with a standard deviation
of one sample to suppress single-bin noise.  The remaining data are fit to the
operational dip estimator
\begin{equation}
    y(\Delta)
    = y_0-A\exp\!\left[-\frac{(\Delta-\Delta_0)^2}{2\sigma^2}\right],
    \label{eq:dip-gaussian-estimator}
\end{equation}
where $A$ is the contrast, $\Delta_0$ is the fitted center, and
\begin{equation}
    \Gamma_{\mathrm{FWHM}}
    =2\sqrt{2\ln2}\,|\sigma|.
\end{equation}
The Gaussian is an operational, consistent linewidth estimator; it is not a
claim that the complete physical response is exactly Gaussian.

Fits are accepted only when all parameters and covariance-derived
uncertainties are finite, $A\ge0.02$,
$\Gamma_{\mathrm{FWHM}}\le0.12~\mathrm{MHz}$, and
$|\Delta_0|\le0.050~\mathrm{MHz}$.  The reported parameter uncertainties are
one standard deviation from the fit covariance matrix.  Of the 100 amplitude
slices, 62 pass these numerical gates.  Thirty are rejected as too broad,
seven have centers outside the accepted window, and one fails both the
contrast and center gates.  Rejected or failed fits remain missing values;
they are never replaced by zero linewidth or zero uncertainty.  The reference
frequency offset for the center plot is $0~\mathrm{MHz}$.
